\chapter{The Permanence Argument}
We now prove that $h_6^2$ survives nontrivially in the sphere Adams
spectral sequence. The calculations leave just one possible nonzero differential,
$d_{12}(h_6^2)=h_1h_4x_{109,12}$. Its existence would require three
relations among classes in stems $124$ through $126$. We show that two
of these relations are incompatible.
As in \cite{LWX126}, the spectra in the quoted statements are
$2$-completed, connective and of finite type. The
\href{paper/2412.10879.pdf}{article PDF} and
\href{paper/main.tex.html}{TeX source} give the original statements.

The Burklund--Xu criterion of Theorem~\ref{sem:Literature:bx-survival}
turns the survival question into a vanishing problem involving a
synthetic theta-five class. We compare the possible choices of that
class, then use the extension rules of Chapter~4 and the calculations
of Chapters~3 and~5 to constrain the remaining obstruction.
The finite-page calculations and the two staircase classes used below
are the inputs of
Theorems~\ref{sem:SSComputation:admitted-finite-inputs}
and~\ref{sem:SSComputation:staircase-inputs}.

The passage from page labels to homotopy classes uses the convergence
and comparison maps of Chapter~2, the finite-quotient comparisons of
Proposition~\ref{sem:Literature:finite-quotient-labels}, and the complete
lifts of Theorem~\ref{sem:Literature:complete-permanent-lifts}.
All deductions in this chapter assume these comparisons and the
computational inputs, with the hypotheses stated there.

The proof proceeds through the survival criterion
(Theorem~\ref{paper:7.3}), the last-differential alternatives
(Proposition~\ref{paper:7.8}), the sixth-quotient representative construction
and the permanence theorem \ref{paper:1.4}. Proposition~\ref{paper:7.9}
is also deduced from the same representative. The finite-page bounds supply
the representatives used in this argument. The final section
records the conditional Toda consequences separately.

\section{The sphere and its obstruction classes}
Let $E_r^{s,t}$ be the sphere Adams spectral sequence, with stem $t-s$,
and let $S^{0,0}$ be the synthetic sphere. Write $F^s$ for the
synthetic Adams filtration. The three recurring second-page labels are
\[
 \begin{aligned}
 W&=x_{126,8,4}+x_{126,8}&&\in E_2^{8,134},\\
 U&=h_0^2x_{124,8}&&\in E_2^{10,134},\\
 T&=h_1h_4x_{109,12}&&\in E_2^{14,139}.
 \end{aligned}
\]
A label detects a homotopy class when their images agree in the
corresponding associated grade. Nonzero detection means that this grade
is nonzero. In particular, nonzero survival of $h_6^2$ requires a common
permanent representative with nonzero image on $E_\infty$.

Write $Q_n=S^{0,0}/\lambda^n$ and $q_n:S^{0,0}\to Q_n$ for quotient
projection. For $k\ge0$, a classical label $x\in E_2^{s,t}$ determines the synthetic
label $\lambda^kx$ in tridegree $(s,t,t-k)$. Detection by this label
in $Q_n$ uses its image under the quotient page map; detection in the
full sphere uses the untruncated comparison. In homotopy products,
$[h_0]\in\pi_{0,1}S^{0,0}$ and $[h_2]\in\pi_{3,4}S^{0,0}$ denote the
chosen lifts of the degree-two and Hopf maps. Their brackets distinguish
them from the corresponding Adams labels. When $h_0$ occurs in a
homotopy product without brackets, it abbreviates $[h_0]$.

\begin{definition}[Admissible representatives and the detection conditions]
\label{sem:Main:admissible-representatives}
\level{coarse}
From \cite{LWX126}, Notation 7.2 and Proposition 7.8.\par

An element $\theta\in\pi_{62,64}S^{0,0}$ is admissible if it is nontrivially detected by $h_5^2$. An element $u\in\pi_{124,134}S^{0,0}$ is admissible if it is nontrivially detected by $U$. In both cases, detection uses the comparison and convergence maps of the synthetic sphere. Representatives may differ by a class of higher filtration.

The possible last differential is
\[
 d_{12}(h_6^2)=T\ne0,\qquad (2,128)\longrightarrow(14,139).
\]

This equation includes the existence of twelfth-page representatives
with nonzero target. Three relations will govern whether it occurs.
The first is
\[
 d_6(W)=0,\qquad (8,134)\longrightarrow(14,139).
\]

This relation includes the assertion that $W$ reaches $E_6$.

The second asks for an admissible $\theta$ whose square
$\theta^2\in\pi_{124,128}S^{0,0}$ is nontrivially detected by
$\lambda^6U$, of tridegree $(10,134,128)$. We call this nonzero
theta-square detection at $\theta$, and ask whether it holds for some
admissible choice.

The third asks for an admissible $u$ such that
$\lambda^3\eta u\in\pi_{125,133}S^{0,0}$ is detected by
$\lambda^6T$, of tridegree $(14,139,133)$. We call this eta-extension
detection at $u$, and again ask whether it holds for some admissible
choice. Here detection allows a zero associated graded class.
We write $\theta_5$ for an admissible theta choice when none has been
specified; the superscript in $\theta_5^{\mathrm{BX}}$ denotes
the distinguished Burklund--Xu choice.
\lean{KIP126.Main.ThetaChoice,KIP126.Main.UChoice,KIP126.Main.H6SquareSupportsD12,KIP126.Main.Stem126D6Vanishes,KIP126.Main.Theta5SquareNonzeroDetectionAt,KIP126.Main.Theta5SquareNonzeroDetection,KIP126.Main.EtaExtensionDetectionAt,KIP126.Main.EtaExtensionDetection}
\end{definition}

\section{Choice independence and the inductive criterion}

\begin{remark}[Choice independence and the total boundary]
\label{sem:Main:remark-7.5}
\level{coarse}
From \cite{LWX126}, Remark 7.5.\par
For every admissible $\theta\in\pi_{62,64}S^{0,0}$,
\[
 \lambda\eta\theta^2=\lambda\eta(\theta_5^{\mathrm{BX}})^2.
\]
Thus the obstruction $\lambda\eta\theta^2$ is independent of the
admissible choice, even though the class $\theta$ may vary.

\href{paper/main.tex.html#L2133}{Original remark}.

For every admissible $\theta\in\pi_{62,64}S^{0,0}$, the total boundary of the first-quotient class $h_6^2$ satisfies
\[
 \delta_1(h_6^2)=\lambda\eta\theta^2.
\]
Here $\delta_1$ is the connecting map for the first lambda quotient.
\uses{sem:Main:admissible-representatives,sem:Literature:bx-total-boundary,sem:Literature:theta-five}
\lean{KIP126.Main.lambda_eta_theta_square_independent,KIP126.Main.delta_h6_square_eq_lambda_eta_theta_square}
\end{remark}
\begin{proof}
\uses{sem:Computation:table-certification,sem:Computation:higher-page-vanishing,sem:Computation:correction-survival,sem:Computation:products-vanishing,paper:7.6,sem:Literature:synthetic-h-zero,sem:Literature:filtered-product-detection,sem:Literature:tmf-hurewicz,sem:Literature:realization-kernel,sem:Literature:complete-detection,sem:AdamsSynthetic:finite-separation}
Put $\theta_0=\theta_5^{\mathrm{BX}}$ and $d=\theta-\theta_0$.
The two choices have the same leading label, so $d\in F^3$.
The finite-page calculations
$E_3^{3,65}=E_3^{4,66}=E_4^{5,67}=0$ raise this to $d\in F^6$.
Since $2\theta_0=0$, graded commutativity gives
\[
 b=\theta^2-\theta_0^2=d^2\in F^{12}\pi_{124,128}S^{0,0}.
\]

Here is the high-filtration calculation used on $b$. The vanishings
$E_5^{11,135}=0$ and $E_4^{12,136}=0$, together with the
filtration-thirteen candidate $C=e_0\Delta h_6g$, show that an element
of $F^{11}\pi_{124,128}$ is either detected by $\lambda^9C$ or is
divisible by $\lambda^{10}$. The relation $h_1C=0$, product detection
and the filtration--$\lambda$ comparison put its $\eta$ multiple in
$F^{15}$. Consequently $a=\eta b\in F^{15}\pi_{125,130}$.
The relations $\lambda[h_0]=2$ and $[h_0]\eta=0$ also give $2a=0$.

The tmf calculation controls this last group as follows:
an element $a\in F^{15}\pi_{125,130}$ with $2a=0$ and zero synthetic
tmf image is zero. Indeed the page calculations eliminate the grades
from fifteen through twenty-four. The only possible grade at twenty-five
is represented by $H=g^4\Delta h_1g$, whose realization has the
tmf image of Proposition~\ref{sem:Literature:tmf-hurewicz}, which is
not annihilated by any positive odd integer. If that grade is nonzero,
subtract $\lambda^{20}\widetilde v$, where $\widetilde v$ is a lift of
$v=\bar\kappa^4w$ detected by $H$. The difference lies in $F^{26}$;
its realization is odd torsion by sphere separation. Multiplying by
an odd integer annihilating this difference contradicts the stated
tmf property, since $a$ has zero detector image. Thus $a\in F^{26}$
is infinitely filtered. Its realization is both odd torsion and killed
by two, hence zero. The realization-kernel comparison makes $a$
$\lambda$-torsion, and finite-quotient separation then gives $a=0$.
Finally, the stem-$62$ tmf vanishing, the support condition
$E_2^{s,63+s}(\mathit{tmf})=0$ for $s\le0$, and the realization and
$\lambda$-kernel comparisons make the synthetic detector image of
each $\theta\in\pi_{62,64}$ zero. Hence the same holds for
$a=\eta(\theta^2-\theta_0^2)$, so this criterion applies.
Multiplying by $\lambda$ proves choice independence.

The total-boundary formula at $\theta_0$ is part of the comparison in
Definition~\ref{sem:Literature:bx-total-boundary}. Substitute the equality
just proved to obtain the formula for every admissible $\theta$.
\end{proof}

% AimPaper: 7.3 theorem 2117 2125
\begin{theorem}[A synthetic criterion for the survival of $h_6^2$]
\label{paper:7.3}
\label{sem:Main:any-choice-criterion}
\label{thm:bjmbx}
\source{AimPaper Theorem 7.3}
\level{coarse}
The Barratt--Jones--Mahowald criterion, in the synthetic form of
\cite[Proposition 7.19]{BurklundXu}, is expressed in the present grading
by \cite[Theorem 7.3 and Remarks 7.4--7.5]{LWX126}
(\href{paper/main.tex.html#L2117}{original statement}).

For every admissible $\theta\in\pi_{62,64}S^{0,0}$ and every integer
$r\geq0$,
\[
 h_6^2\text{ survives to the classical Adams }E_{r+3}\text{-page}
 \quad\Longleftrightarrow\quad
 [\lambda\eta\theta^2]=0\text{ in }\pi_{125,129}(S^{0,0}/\lambda^{r+1}).
\]
Moreover,
\[
 h_6^2\text{ survives nontrivially to }E_\infty
 \quad\Longleftrightarrow\quad\lambda\eta\theta^2=0
 \text{ in }\pi_{125,129}S^{0,0}.
\]
Admissible choices of $\theta$ exist by
Proposition~\ref{sem:Literature:theta-five}. Thus each equivalence also
holds with ``for some admissible $\theta$'' on the right, as in the
original theorem.
\uses{sem:Main:admissible-representatives,paper:3.3,sem:Literature:theta-five,sem:Literature:hopf-classes,sem:Literature:bx-choice}
\lean{KIP126.Main.any_choice_criterion}
\end{theorem}
\begin{proof}
\uses{sem:Literature:bx-survival,sem:Literature:bx-total-boundary,sem:Main:remark-7.5}
First compare the two forms of the quotient condition, as in
\cite[Remark 7.4]{LWX126}. For the quotient projection $q_n$ and
$a\in\pi_{125,130}S^{0,0}$, one has, for every $n\geq0$,
\[
 q_n(a)=0\quad\Longleftrightarrow\quad q_{n+1}(\lambda a)=0.
\]
In particular, with $a=\eta(\theta_5^{\mathrm{BX}})^2$ this identifies the two quotient conditions
\[
 q_r\bigl(\eta(\theta_5^{\mathrm{BX}})^2\bigr)=0
 \quad\Longleftrightarrow\quad
 q_{r+1}\bigl(\lambda\eta(\theta_5^{\mathrm{BX}})^2\bigr)=0.
\]

Exactness identifies $\ker q_n$ with the image of $\lambda^n$. In this degree multiplication by $\lambda$ is injective, so $\lambda a=\lambda^{n+1}b$ is equivalent to $a=\lambda^n b$. This proves the quotient comparison used in Remark 7.4 in the stated degree.

For $r\geq1$, apply the distinguished-choice criterion of
Theorem~\ref{sem:Literature:bx-survival} and this quotient comparison.
Remark~\ref{sem:Main:remark-7.5} gives
\[
 \lambda\eta\theta^2=\lambda\eta(\theta_5^{\mathrm{BX}})^2
\]
for every admissible $\theta$, so the criterion holds for every such
choice. At $r=0$, $h_6^2$ survives to $E_3$ independently, and every
lambda multiple maps to zero in the first quotient. The untruncated
equivalence follows from the distinguished-choice criterion and the
same choice independence. Existence of an admissible choice then gives
the original existential formulations.
\end{proof}

% AimPaper: 7.10 lemma 2221 2228
\begin{lemma}[Independence of the theta choice]
\label{paper:7.10}
\source{AimPaper Lemma 7.10}
\level{coarse}
From \cite{LWX126}, Lemma 7.10.\par

\label{lem:equistate4}
The theta-square nonzero-detection condition of
Definition~\ref{sem:Main:admissible-representatives} holds for some
admissible $\theta$ if and only if it holds for every admissible
$\theta$. In that case, for each choice there is an admissible
$u\in\pi_{124,134}S^{0,0}$ with
\[
 \theta^2=\lambda^6u\ne0\quad\text{in }\pi_{124,128}S^{0,0}.
\]

\uses{sem:Main:admissible-representatives,sem:Literature:theta-five,sem:Literature:complete-detection,sem:Literature:filtered-product-detection,paper:3.11}

\lean{KIP126.Main.theta5_square_nonzero_detection_choice_equivalence}

\end{lemma}
\begin{proof}
\uses{sem:Literature:complete-boundary-lift}
\textit{(See \cite{LWX126}, Lemma 7.10.)} \href{paper/main.tex.html#L2230}{Original proof}.

Only the implication from one admissible choice to every admissible choice needs proof. The two-primary stem-62 computation of \cite{IWX} places the difference between two classes with the same $h_5^2$ detection in Adams filtration at least six. The realization and filtration comparisons give the same estimate for synthetic choices in this degree. The exponent-two square identity then puts the difference of their squares in filtration at least twelve.

The detector $\lambda^6h_0^2x_{124,8}$ has filtration ten. Adding a term
in filtration at least twelve preserves this nonzero leading term.
Thus detection for one choice implies the same nonzero detection for
every choice. The exponent-two calculation applies to the two-primary
subgroup of the ordinary $62$-stem.
The filtration--$\lambda$ comparison supplies a weight-$134$ preimage
of $\theta^2$ under $\lambda^6$. The window comparison determines its
label modulo $B_7$; the torsion lifts of those boundary labels correct
it to $U$ without changing its $\lambda^6$ image. This gives the
displayed equality with an admissible $u$.
\end{proof}

% AimPaper: 7.11 lemma 2244 2252
\begin{lemma}[Independence of the detected $U$ representative]
\label{paper:7.11}
\source{AimPaper Lemma 7.11}
\level{coarse}
From \cite{LWX126}, Lemma 7.11.\par

\label{lem:equistate5}
The eta-extension detection condition of
Definition~\ref{sem:Main:admissible-representatives} holds for some
admissible $u$ if and only if it holds for every admissible $u$.
In that case each such $u$ admits a class
$t_T\in\pi_{125,139}S^{0,0}$ detected by $T$ such that
\[
 \lambda^3\eta u=\lambda^6t_T\quad\text{in }\pi_{125,133}S^{0,0}.
\]

\uses{sem:Main:admissible-representatives,paper:3.6,paper:3.12,sem:Computation:table-certification,sem:Computation:products-vanishing,sem:Computation:correction-survival,sem:Computation:higher-page-vanishing,sem:Computation:low-differentials,sem:Computation:literature-labels,sem:Literature:filtered-product-detection,sem:Literature:complete-detection,sem:Literature:tmf-hurewicz,sem:Literature:sphere-vanishing,sem:AdamsSynthetic:finite-separation}

\lean{KIP126.Main.eta_extension_detection_choice_equivalence}

\end{lemma}
\begin{proof}
\uses{sem:Literature:complete-boundary-lift}
\textit{(See \cite{LWX126}, Lemma 7.11.)} \href{paper/main.tex.html#L2254}{Original proof}.

It suffices to show that the difference of two admissible $U$
representatives, multiplied by $\lambda^3\eta$, lies in $F^{15}$.
Such a change cannot affect the target grade at filtration fourteen.

The label $U=h_0^2x_{124,8}$ has filtration ten, so two admissible
representatives differ in filtration at least eleven in stem $124$.
The $d_2$ and $d_3$ calculations eliminate the leading terms in
filtrations eleven and twelve. Every possible filtration-thirteen
term is annihilated by $h_1$ in Ext. Multiplication by $\lambda^3\eta$
therefore puts the difference in filtration at least fifteen.
The filtration--$\lambda$ comparison gives a weight-$139$ preimage of
$\lambda^3\eta u$ under $\lambda^6$. As in the preceding proof, its
$B_7$ label ambiguity can be corrected by torsion lifts killed by
$\lambda^6$, giving $t_T$ with label $T$ and the displayed equality.
\end{proof}

\section{The last possible differential}

% AimPaper: 7.8 proposition 2181 2204
\begin{proposition}[The last differential and its detection conditions]
\label{paper:7.8}
\source{AimPaper Proposition 7.8}
\level{coarse}
From \cite{LWX126}, Proposition 7.8.\par

\label{prop:possible_h_6_sq}\label{prop:possibleh62}
Exactly one of the following two statements is true:
\begin{enumerate}
    \item The element $h_6^2$ survives nontrivially to the $E_\infty$-page.

    \item There is a nonzero classical Adams differential
    $$d_{12}(h_6^2) = h_1h_4x_{109,12}.$$
\end{enumerate}
Furthermore, statement $(2)$ is true if and only if the following three statements are all true:
\begin{enumerate} \setcounter{enumi}{2}
    \item The class $W$ reaches $E_6$ and satisfies $d_6(W)=0$.

    \item For some admissible $\theta_5$, the square $\theta_5^2$ is
    nontrivially detected by $\lambda^6U$. In that case an admissible
    $u\in\pi_{124,134}S^{0,0}$ can be chosen such that
    \[
      \theta_5^2=\lambda^6u\ne0\quad\text{in }\pi_{124,128}S^{0,0}.
    \]

    \item For some admissible $u$, the class $\lambda^3\eta u$ is
    detected by $\lambda^6T$. In that case there is a class
    $t_T\in\pi_{125,139}S^{0,0}$ detected by $T$ with
    \[
      \lambda^3\eta u=\lambda^6t_T
        \quad\text{in }\pi_{125,133}S^{0,0}.
    \]
\end{enumerate}

\uses{paper:7.3,paper:7.6,paper:7.10,paper:7.11,paper:3.6,sem:Literature:finite-bockstein,paper:4.10,sem:Literature:tmf-hurewicz,sem:Literature:bx-total-boundary,sem:Computation:table-certification,sem:Computation:incoming-control,sem:Main:admissible-representatives}

\lean{KIP126.Main.h6_square_differential_dichotomy_and_criterion}

\end{proposition}
\begin{proof}
\uses{sem:Main:admissible-representatives,sem:Computation:w-t-u,sem:Computation:correction-survival,sem:Computation:table-bases,sem:Literature:complete-permanent-lifts}
The stem-$124$ calculation first gives the following alternatives.
For every $\theta\in\pi_{62,64}S^{0,0}$, at least one of the following holds:
\begin{enumerate}
\item $\theta^2=\lambda^6u$ for an admissible $u\in\pi_{124,134}S^{0,0}$;
\item $\theta^2=\lambda^9c$ for some $c\in\pi_{124,137}S^{0,0}$ nontrivially detected by the correction class $e_0\Delta h_6g$ in bidegree $(13,137)$;
\item $\theta^2=\lambda^{10}a$ for some $a\in\pi_{124,138}S^{0,0}$.
\end{enumerate}

\textit{(See \cite{LWX126}, Proposition 7.8.)} \href{paper/main.tex.html#L2300}{Original proof}.

Use the total-boundary formula from the quadratic construction,
\[
 \delta_1(h_6^2)=\lambda\eta\theta_5^2.
\]
The rigidity and boundary-extension comparisons identify the first nonzero detector of this boundary with the first nonzero classical differential on $h_6^2$. The sphere computation leaves only $d_{12}(h_6^2)=h_1h_4x_{109,12}$ as a possible nonzero outcome. The finite-quotient comparison and Lemmas~\ref{paper:7.10} and~\ref{paper:7.11} allow the detection relations to be evaluated at compatible choices.

If sixth-differential vanishing, theta-square nonzero detection and eta-extension detection all hold, multiplication of the two detection relations gives
\[
 \lambda\eta\theta_5^2=\lambda^{10}[h_1h_4x_{109,12}].
\]
By Fact~\ref{paper:7.6}, sixth-differential vanishing excludes the shorter incoming $d_6$ to this target. The boundary is therefore nonzero and gives exactly the stated $d_{12}$.

Conversely the stem-124 computation gives the three possibilities displayed above: detection of $\theta_5^2$ by $\lambda^6h_0^2x_{124,8}$, the filtration-13 correction, or divisibility by $\lambda^{10}$. In the latter two cases the $h_1$ product computation places $\eta\theta_5^2$ in filtration at least fifteen. The remaining class in this range is detected by tmf, whereas $\theta_5$ has zero tmf image, so the product is zero. In the first case failure of eta-extension detection gives the same high-filtration alternative and the same vanishing conclusion. If sixth-differential vanishing fails, the $d_6$ kills the proposed target and the same argument again forces the total boundary to vanish.

Thus failure of any of these three conditions yields permanence by Theorem~\ref{paper:7.3}; their conjunction yields the nonzero $d_{12}$. These alternatives are exclusive and exhaustive.
\end{proof}

% AimPaper: 7.15 fact 2466 2468
\begin{fact}[Nonzero fifth-page survival and incoming control for $Y$]
\label{paper:7.15}
\source{AimPaper Fact 7.15}
\level{coarse}
From \cite{LWX126}, Fact 7.15.\par

\label{fact:h_0_sq_mul_x_125_9_2}\label{fact:h02x1259}
Put $Y=h_0^2x_{125,9,2}\in E_2^{11,136}$. This class survives
nontrivially to $E_5$ and is not the nonzero target of any classical
incoming differential. The assertion concerns its fifth-page class and incoming differentials.

\uses{sem:Computation:w-t-u,sem:Computation:y-to-e5}
\lean{KIP126.Computation.Standard.y_not_hit}
\end{fact}
\begin{proof}
\uses{sem:Computation:y-to-e5,sem:Computation:table-bases,sem:Computation:table-certification,paper:7.8}
Nonzero survival to $E_5$ is the preceding computation. An incoming $d_r$ to $Y=h_0^2x_{125,9,2}$ has source bidegree $(11-r,137-r)$. The finite calculations exclude $r=2,\ldots,6$. At $r=7$ the source has already vanished by the recorded $d_2$ and $d_3$; at $r=8$ its outgoing differential is zero by the recorded $d_2$. The $E_2$ sources for $r=10,11$ are empty, and $r\ge12$ gives negative filtration. The only remaining possibility is
\[
 d_9(h_6^2)=h_0^2x_{125,9,2}.
\]
To exclude this $d_9$, use the last-differential alternative rather
than a longer survival bound on $Y$.

Proposition~\ref{paper:7.8} says that $h_6^2$ is either permanent or
supports the nonzero $d_{12}$ and hence survives to $E_{12}$. Both
alternatives exclude a nonzero $d_9$, proving the no-incoming
assertion.
\end{proof}

\section{Local constructions in the ninth lambda quotient}

We now construct the representatives needed for the obstruction
argument. The first lemma produces a class $\alpha_1$ and correction
terms in $Q_9$; the second computes a $[h_2]$ product using the Hopf
cofiber. Both constructions use the finite-window comparisons of
Chapter~2 and the calculations of Chapters~3 and~5.

% AimPaper: 7.20 lemma 2617 2620
\begin{lemma}[The normalized Hopf-element extension]
\label{paper:7.20}
\source{AimPaper Lemma 7.20}
\level{coarse}
From \cite{LWX126}, Lemma 7.20.\par

\label{lem:nuext125}
    There exists a homotopy class $[\lambda^4 h_1 x_{121,7}] \in \pi_{122,122+4}S^{0,0}/\lambda^9$, such that
    $$[\lambda^4 h_1 x_{121,7}] \cdot [h_2] = \lambda [\lambda^5 h_0^2 x_{125,9,2}] \in \pi_{125,125+5} S^{0,0}/\lambda^9.$$

\uses{paper:7.19,paper:4.6,paper:4.2,sem:Computation:nu-cofiber,sem:Computation:cofiber-results,sem:Computation:top-cell-comparison,sem:Computation:products-vanishing,sem:Computation:y-to-e5,sem:Computation:table-bases,sem:Computation:table-certification,sem:ProofTools:nu-zero-layer,sem:AdamsComputation:direct-d2-records}

\lean{KIP126.Main.Paper.exists_nu_extension_in_ninth_quotient}

\end{lemma}
\begin{proof}
\textit{(See \cite{LWX126}, Lemma 7.20.)} \href{paper/main.tex.html#L2622}{Original proof}.

Fact~\ref{paper:7.19} and the finite-quotient comparison supply a representative detected by $\lambda^4h_1x_{121,7}$ in the length-nine quotient. To compute its $h_2$ image, use the cofiber triangle $S^3\xrightarrow{\nu}S^0\to S^0/\nu\to S^4$ and the top-cell comparison.

The cofiber calculation gives a class
\[
 \bar x=h_1x_{121,7}[4]+x_{126,8}[0]+x_{126,8,2}[0]
\]
whose top-cell image is $h_1x_{121,7}$ and whose $d_3$ is
$h_0^2x_{125,9,2}[0]$. The bottom-cell inclusion detects
$h_0^2x_{125,9,2}$; its $h_2$ nondivisibility prevents it from
disappearing under this inclusion. Apply the zero-layer $\nu$ comparison
of Theorem~\ref{sem:ProofTools:nu-zero-layer} to this class and its
$d_3$. This gives the required finite $h_2$-extension. The argument uses this specific zero-layer comparison.

For the target detector, put $Y=h_0^2x_{125,9,2}$. Its nonzero
fifth-page representative from Fact~\ref{sem:Computation:y-to-e5}
excludes incoming differentials on pages $2$ through $4$. The complete
bases and finite-page calculations give $E_4^{6,132}=0$ and $E_6^{5,131}=0$;
these exclude the sources of incoming $d_5$ and $d_6$, respectively.
Thus no differential on a page $2\leq r\leq6$ hits $Y$. This finite
window supplies the nonzero target detector in the ninth quotient.

Lift the relation from the length-three to the length-five quotient. The only possible additional filtration-ten correction is excluded by the product and finite-page calculations. Push the resulting relation along $\lambda^4: \Sigma^{0,-4}S^{0,0}/\lambda^5\to S^{0,0}/\lambda^9$. Naturality of the multiplication gives exactly
\[
 [\lambda^4h_1x_{121,7}][h_2]
       =\lambda[\lambda^5h_0^2x_{125,9,2}],
\]
with representatives in the quotients specified in the statement.
\end{proof}

\section{The sixth-quotient representative}

\begin{lemma}[A sixth-quotient representative and a permanence criterion]
\label{sem:Main:sixth-quotient-representative}
\level{coarse}
Work in the fixed sphere and synthetic comparison system.
Assume the convergence isomorphisms, the comparison of the degree-two
map with $h_0$, and the finite-page computations used above.
Let $C_2$ and
$\bar x\in E_2^{8,132}(C_2)$ be the completed Moore spectrum and
bottom-cell class of Definition~\ref{sem:AdamsComputation:moore-object}. Put
\[
 V=x_{123,9}+h_0x_{123,8}\in E_2^{9,132}(S^0),\qquad
 Q_n=S^{0,0}/\lambda^n,
\]
and write $F^s$ for the synthetic Adams filtration. Assume the following
two survival conditions:
\[
 \bar x\in Z_4^{8,132}(C_2),\qquad
 V\in Z_8^{9,132}(S^0).
\]
These conditions say that the indicated classes reach $E_5$ and
$E_9$, respectively. Then there is a class
\[
 \beta_6\in\pi_{123,129}Q_6
\]
detected by $\lambda^3V$ and satisfying $h_0\beta_6=0$.
Consequently, $h_6^2\in E_2^{2,128}(S^0)$ has a common permanent representative
whose $E_\infty$ image is nonzero.
Theorem~\ref{sem:SSComputation:selected-class-survival} supplies these two survival
conditions from the ordinary computation inputs of
\ref{sem:SSComputation:ordinary-differential-inputs}, which follow from
\ref{sem:SSComputation:admitted-finite-inputs}.

\lean{KIP126.Main.Internal.Moore.moore_sixth_representative_of_inputs,KIP126.Main.Internal.Moore.h6_sq_permanent_of_moore_inputs}

\uses{sem:Main:admissible-representatives,sem:AdamsComputation:moore-object,sem:AdamsSynthetic:lambda-quotients,sem:Literature:finite-quotient-labels}
\end{lemma}
\begin{proof}
\uses{paper:7.8,paper:7.11,paper:7.20,paper:7.21,sem:SpectralSequence:pairing,sem:AdamsComputation:moore-scope,sem:SSComputation:selected-class-survival,sem:Computation:table-certification,sem:Computation:higher-page-vanishing,sem:Computation:cofiber-results,sem:Computation:product-boundaries,sem:Literature:low-toda-comparison,sem:Literature:symmetric-toda-comparison,sem:SSComputation:staircase-inputs,sem:Literature:sphere-vanishing,sem:Literature:complete-permanent-lifts,sem:Literature:infinite-finite-window,sem:Literature:infinite-restriction,sem:Literature:normalized-e2-long-exact,sem:Literature:normalized-finite-quotient-row,sem:Literature:quotient-algebras,sem:Literature:first-quotient-products,sem:Literature:filtered-product-detection,sem:Literature:may-tensor,sem:Literature:toda-functoriality,sem:ProofTools:finite-strip,sem:AdamsComputation:direct-d2-records}
First, the sphere calculation implies $\bar x\ne0$. Write
$x=b_{S;8,132,0}$, so that $\bar x=i_*x$. If $i_*x=0$, the
normalized $E_2$ exact sequence puts $x$ in the image of the operation
induced by the normalized doubling map. Its first-quotient comparison
shows that every class in this image is a multiple of $h_0$.
Thus $x=h_0z$ for some $z\in E_2^{7,131}(S^0)$. The sphere
calculation gives $d_2z=0$ for every class in this source degree, and
$d_2h_0=0$. Leibniz then gives $d_2x=0$, contrary to the recorded
nonzero $d_2$ on $x$. Hence $\bar x\ne0$.

Now construct the sixth-quotient representative. The finite-window
comparison and $\bar x\in Z_4$ give
$c_4\in\pi_{124,132}(\nu C_2/\lambda^4)$ detected by $\bar x$.
Take its top-cell image and use the suspension isomorphism
to obtain $a_0\in\pi_{123,131}Q_4$. Since $q_*\bar x=0$, this image
lies in $F^9$.

It cannot lie in $F^{10}$. Otherwise the top-cell image of $c_4$,
before desuspension, would also lie in $F^{10}$. Its restriction to the
second quotient would vanish, by the finite high-filtration bound. Exactness of
the normalized Moore quotient row would then lift the restriction
$c_2$ of $c_4$ to the sphere second quotient. The specified sphere
calculation $Z_2^{8,132}(S^0)=0$ puts every such lift in $F^9$, hence
also $c_2\in F^9$. But $c_2$ is still detected by $\bar x\ne0$; its
finite-window denominator is $B_1=0$, so its filtration-eight grade is
nonzero. This is a contradiction.

In the normalized Moore row, the ordinary suspended top image is killed
by the quotient map induced by $h_0$. The commutative algebra structure
on $Q_4$ identifies this map, under the suspension and unit isomorphisms, with the
$h_0$ action. Tensoring with a bigraded sphere preserves and reflects the
Adams filtration. Transport through these comparisons therefore
gives a representative $a\in\pi_{123,131}Q_4$ with
$a\in F^9\setminus F^{10}$ and $h_0a=0$. The sphere calculation
$Z_3^{9,132}(S^0)=\{0,V\}$ identifies its detector as $\lambda V$.
The representatives $a_0$ and $a$ have the same detector.

The layer map $\Sigma^{0,-2}Q_4\to Q_6$ induced by $\lambda^2$ sends
$a$ to $\beta_6$. Its detection comparison applies because
$V\in Z_8\subseteq Z_3$, and gives the detector $\lambda^3V$.
Compatibility with the sphere action gives $h_0\beta_6=0$.

Using this representative and $V\in Z_8$, the finite-window comparison
gives a class
$\beta_{11}\in\pi_{123,129}Q_{11}$ detected by $\lambda^3V$.
Its restriction to $Q_6$ has the same detector as $\beta_6$, so their
difference lies in $F^{10}$. The recorded low-page calculation
$Z_2/B_5(10,133)=0$ raises this difference to $F^{11}$; multiplication
by $h_0$ raises it to $F^{12}$. Since $h_0\beta_6=0$, the restriction
of $h_0\beta_{11}$ has filtration at least twelve.

The product $h_0\beta_{11}$ already lies in $F^{11}$. On its possible
leading grade, restriction from $Q_{11}$ to $Q_6$ is the injection
\[
 Z_7/B_5(11,134)\longrightarrow Z_2/B_5(11,134).
\]
The denominator is the same on both sides. Therefore
$h_0\beta_{11}\in F^{12}$. The higher finite-window vanishings raise
this to $F^{16}$, and restriction to $Q_9$ gives $h_0\beta_9=0$,
where $\beta_9$ is the restriction of this same $\beta_{11}$.

If the nonzero $d_{12}$ alternative in Proposition~\ref{paper:7.8}
held, it would supply sixth-differential vanishing and a compatible
representative $u$ satisfying the eta-extension condition. The
stem-124 correction calculation supplies classes $a_2,a_3$ in $Q_9$
such that
\[
 \eta\beta_9=\lambda^3q_9(u)+\lambda^6a_2,
 \qquad \eta a_2=\lambda a_3.
\]
Its filtration-fourteen and filtration-fifteen corrections use the
two staircase inputs of \ref{sem:SSComputation:staircase-inputs},
the sphere vanishing line and the permanent-lift comparison.
The relation $h_0\beta_9=0$ permits the required Toda construction.
The target and representative calculations then force a differential of
length at most five hitting $h_1h_4x_{109,12}[0]$ in the $\nu$ cofiber,
contradicting its no-incoming window. Proposition~\ref{paper:7.8} therefore
gives the required nonzero permanence.

Finally instantiate the two survival premises with
Theorem~\ref{sem:SSComputation:selected-class-survival}. This last step uses the
two ordinary computation inputs stated in Chapter~5.
\end{proof}

\section{Excluding the eta-extension obstruction}

% AimPaper: 7.9 proposition 2208 2210
\begin{proposition}[The sixth-differential vanishing excludes the eta extension]
\label{paper:7.9}
\source{AimPaper Proposition 7.9}
\level{coarse}
From \cite{LWX126}, Proposition 7.9.\par

\label{prop:state5false}
If $d_6(W)=0$, the eta-extension detection condition of
Definition~\ref{sem:Main:admissible-representatives} fails for every
admissible $u$. Equivalently, the sixth-differential vanishing and
eta-extension detection conditions cannot hold simultaneously.

\uses{sem:Main:sixth-quotient-representative,sem:Computation:cofiber-results,paper:7.11,paper:7.20,paper:7.21}

\lean{KIP126.Main.eta_extension_excluded_of_d6_vanishes}

\end{proposition}
\begin{proof}
\uses{sem:Main:sixth-quotient-representative,sem:SSComputation:selected-class-survival,paper:7.11,paper:7.20,paper:7.21,sem:Literature:low-toda-comparison,sem:Literature:symmetric-toda-comparison,sem:Literature:infinite-finite-window,sem:Literature:may-tensor,sem:Computation:cofiber-results,sem:SSComputation:staircase-inputs}
Use the sixth-quotient construction to obtain a representative
$\beta_6$ detected by $\lambda^3V$ and killed by $h_0$.
The $E_9$ continuation of $V$ supplies a compatible length-eleven
representative. Comparing its restriction with $\beta_6$ and using the
finite-window gaps proves that its ninth-quotient image $\beta_9$
also satisfies $h_0\beta_9=0$.

Suppose sixth-differential vanishing and eta-extension detection hold,
and choose the admissible $u$ occurring in the latter condition.
The correction calculation, including the two staircase inputs, gives
\[
 \eta\beta_9=\lambda^3q_9(u)+\lambda^6a_2,
 \qquad \eta a_2=\lambda a_3.
\]
The vanishing $h_0\beta_9=0$ defines the required Toda bracket.
The low and symmetric Toda comparisons, the ninth-quotient Hopf extension
and the multiplied target argument force an incoming differential of
length at most five on $T[0]$ in the $\nu$ cofiber. The cofiber calculation
excludes this.
\end{proof}

\section{Permanence and the Adams second-line classification}

% AimPaper: 7.1 theorem 2106 2108
% AimPaper: 1.4 theorem 128 130
\begin{theorem}[Permanence of $h_6^2$]
\label{paper:1.4}
\label{paper:7.1}
\label{thm:126survives}
\level{coarse}
From \cite{LWX126}, Theorem 1.4, Theorem 7.1.\par

\label{thm:h_6_sq}\label{thm:h62}
The element $h_6^2$ survives nontrivially to the $E_\infty$-page in the
Adams spectral sequence.

\uses{sem:Main:sixth-quotient-representative,sem:SSComputation:selected-class-survival,sem:AdamsSynthetic:standard-classes}

The conclusion is independent of the normalization of the standard
classes. Let $E$ be the sphere Adams spectral sequence with the
second-page multiplication of \ref{sem:AdamsSynthetic:standard-classes},
and let $H$ be any standard class family in $E$, so
$H_i\in E_2^{1,2^i}$ and $H_i\ne0$ for $0\le i\le7$. Then
\[
H_6\cdot H_6\in E_2^{2,128}
\]
survives nontrivially to $E_\infty$: it has a common permanent
representative whose limiting image is nonzero.

\lean{KIP126.Main.h6_sq_permanent}

\end{theorem}
\begin{proof}
\uses{sem:Main:sixth-quotient-representative,sem:SSComputation:selected-class-survival,sem:Computation:h6-line}
\textit{(See \cite{LWX126}, Theorem 1.4.)} \href{paper/main.tex.html#L2212}{Original proof}.

Apply the sixth-quotient representative construction to the two page bounds of
Theorem~\ref{sem:SSComputation:selected-class-survival}. It constructs the
$Q_6$ witness and excludes the nonzero $d_{12}$ alternative, giving
nonzero permanence.

This proves nonzero permanence for $h_6^2$. For an
arbitrary standard sphere class family $H$, its defining nonzero condition
gives $H_6\ne0$. Theorem~\ref{sem:Computation:h6-line} therefore gives
$H_6=h_6$. Their squares are equal under the same second-page multiplication,
which supplies the stated nonzero permanence for $H_6\cdot H_6$.
\end{proof}

\begin{theorem}[The standard generators are nonzero on the second page]
\label{sem:Main:standard-hi-nonzero}
\level{coarse}
\uses{sem:AdamsSynthetic:standard-classes}
Let $H$ be any standard sphere class family as in
Theorem~\ref{paper:1.4}. For every $0\le i\le7$,
\[
 H_i\in E_2^{1,2^i}(S^0),\qquad H_i\ne0.
\]
\lean{KIP126.Main.standard_hi_ne_zero}
\end{theorem}
\begin{proof}
Nonvanishing is part of the definition of a standard class family.
For the chosen family $h_i=\psi_i(e_i)$, this condition follows from
$e_i\ne0$ and the injectivity of the comparison $\psi_i$ in
Definition~\ref{sem:AdamsSynthetic:hi-data}.
\end{proof}

\begin{theorem}[The square is nonzero on the second page]
\label{sem:Main:h6-square-e2-nonzero}
\level{coarse}
\uses{sem:Main:standard-hi-nonzero,sem:Computation:h6-line,sem:Computation:square-line}
For every standard sphere class family $H$ as in
Theorem~\ref{paper:1.4}, the same product as in that theorem satisfies
\[
 H_6^2=H_6\cdot H_6\in E_2^{2,128}(S^0)\setminus\{0\}.
\]
Here the bidegree convention is $(s,t)$: Adams filtration $2$,
internal degree $128$, and stem $128-2=126$.
\lean{KIP126.Main.h6_square_e2_nonzero}
\end{theorem}
\begin{proof}
By Theorem~\ref{sem:Main:standard-hi-nonzero}, $H_6\ne0$.
Theorem~\ref{sem:Computation:h6-line} identifies the unique nonzero
class of $E_2^{1,64}$ as $h_6$, so $H_6=h_6$.
Theorem~\ref{sem:Computation:square-line} gives $h_6^2\ne0$ already
on the second page. The product has bidegree
$(1,64)+(1,64)=(2,128)$. This argument uses the initial-page
calculation independently of the permanence theorem.
\end{proof}

% AimPaper: 1.5 corollary 151 154
\begin{corollary}[Survivors on the finite Adams second line]
\label{sem:Main:finite-two-line-classification}
\source{AimPaper Corollary 1.5}
\level{coarse}
A finite-range revision of \cite{LWX126}, Corollary 1.5.\par
For every integer $t\leq256$, a class $x\in E_2^{2,t}$ survives
nontrivially to $E_\infty$ if and only if it is one of
\[
 h_0h_2,\quad h_0h_3,\quad h_2h_4,\quad h_1h_j\ (3\leq j\leq7),\quad h_j^2\ (0\leq j\leq6).
\]
Each listed product is understood in its own bidegree on the sphere
second page; an equality to $x$ includes that bidegree identification.
\uses{paper:1.4,sem:E2Computation:two-line-exhaustive,sem:AdamsCalculations:two-line-d2-screening,sem:Literature:two-line-classical-inputs}
\lean{KIP126.Main.adams_two_line_classification}
\end{corollary}
\begin{proof}
Write any nonzero second-line class as a standard product using the exhaustive $E_2$ basis. Its third-page continuation restricts it to the corrected Adams--Lin--Leibniz candidates, including $h_0h_4$. Four exceptional mixed products support nonzero outgoing differentials, while the bounded Browder criterion and geometric HHR input exclude the sole remaining finite square of index seven. Conversely, the remaining products survive by the low-dimensional, Mahowald-family and Kervaire literature inputs, together with the main $h_6^2$ theorem.
\end{proof}

\section{Conditional Toda consequences}

The two conditional conclusions of \cite{LWX126}, Lemma 7.16 and
Corollary 7.18, follow here from Proposition~\ref{paper:7.9}, which
excludes their joint hypotheses.

% AimPaper: 7.16 lemma 2470 2474
% AimPaper: 7.18 corollary 2571 2575
\begin{lemma}[Conditional Toda detection and the $h_0$ extension]
\label{paper:7.16}
\label{lem:toda2ext}
\label{paper:7.18}
\label{cor:2ext125}
\level{coarse}
\source{AimPaper Lemma 7.16}
\source{AimPaper Corollary 7.18}
\uses{paper:7.15,paper:7.6,sem:Main:admissible-representatives,sem:Literature:low-toda-comparison,sem:Literature:symmetric-toda-comparison,sem:Literature:toda-functoriality,sem:Computation:nondivisibility,sem:Computation:products-vanishing,sem:Computation:low-differentials,sem:AdamsComputation:d2-h6-calibration,sem:Computation:table-certification,paper:7.9,sem:SpectralSequence:detection,sem:Literature:filtered-product-detection}
From \cite{LWX126}, Lemma 7.16 and Corollary 7.18.\par
Assume sixth-differential vanishing and eta-extension detection, using
Lemma~\ref{paper:7.11} to choose the admissible $U$ representative.
Let $\alpha_1$ be the class of Lemma~7.14 of \cite{LWX126} and put
\[
 \mathcal T=\langle\lambda^3\alpha_1,[h_0],\eta\rangle
 \subseteq\pi_{125,132}Q_9.
\]
Then $0\notin\mathcal T$, and its members are nontrivially detected by
$\lambda^4Y$, where $Y=h_0^2x_{125,9,2}$. Moreover, for every class
$y\in\pi_{125,132}Q_9$ nontrivially detected by $\lambda^4Y$, there is
a class $t_T\in\pi_{125,139}Q_9$ nontrivially detected by
$T=h_1h_4x_{109,12}$ such that
\[
 y[h_0]=\lambda^6t_T\ne0\quad\text{in }\pi_{125,133}Q_9.
\]
\lean{KIP126.Main.Paper.exists_alpha_one_with_toda_bracket_detected_by_Y,KIP126.Main.Paper.exists_syntheticH0_extension_of_Y_detection}
\end{lemma}
\begin{proof}
Proposition~\ref{paper:7.9} excludes the joint hypotheses, proving
both conditional conclusions. Its independent proof uses the multiplied
relation constructed there. For the corresponding paper arguments, see
\href{paper/main.tex.html#L2482}{the proof of Lemma 7.16} and
\href{paper/main.tex.html#L2577}{the proof of Corollary 7.18}.

For the representative interpretation, distinguish bracket membership
from the multiplied value. The relations of Lemma~7.14 of \cite{LWX126}
define $\mathcal T$. The corrections in Proposition~\ref{paper:7.9} preserve
$\lambda^2[h_0]z$. The representative comparison for the bracket uses a
member in $F^{11}$, control of its restriction to $Q_3$, and membership
of the corrected representative in $\mathcal T$.

Likewise, if $y_0\in\mathcal T$ satisfies
$[h_0]y_0=\lambda^6t_0$, then another representative $y$ with the same
$Y$ detector satisfies exactly
\[
 [h_0]y=\lambda^6t_0+[h_0](y-y_0).
\]
All admissible $T$ representatives have the same $\lambda^6$ image in
this degree. Thus the unmultiplied equation for $y$ is equivalent to
$[h_0](y-y_0)=0$. The permanence argument instead uses the
vanishing of $\lambda^2[h_0](y-y_0)$. The kernel of $\lambda^2$ controls the passage to an unmultiplied
equation; for an $\eta$-multiplied equation, the corresponding comparison
uses the kernel of $\eta$ multiplication.
\end{proof}
